% GENERATED FILE. Edit canonical lessons, then run npm run generate:books.
% canonical-source-hash: fnv1a64:2fa85dc071b1cb03
% canonical-lessons: PT-C364-protestar, PT-C364-protesto, PT-C364-manifestacao, PT-C364-manifestar, PT-C364-exigir

\chapter{Na rua}
\label{ch:pt-na-rua}
\hlchaptermodality{364}

\begin{chapteropening}
\textbf{By the end of this chapter:} \emph{I can talk about protests: what people demand, demonstrations and how they show their anger.}

Complete the last lesson of chapter 364: I can talk about protests: what people demand, demonstrations and how they show their anger.
\end{chapteropening}

\section[protestar]{protestar — to protest}

\label{lesson:PT-C364-protestar}



\begin{quote}
Before the new one: say the Portuguese for the freedom, then the Portuguese for the equality.
\end{quote}

\subsection*{You'll want to know: protestar}
\textbf{protestar} — \textquotedblleft{}to protest\textquotedblright{}.

\begin{quote}
Os moradores protestaram contra o barulho.
\end{quote}

\begin{quote}
\emph{The residents protested against the noise.}
\end{quote}

\emph{Protestar contra} — to protest against.

\begin{cousinweb}
\emph{protestar} is a learned borrowing from Latin \emph{prōtestārī} \textquotedblleft{}to declare in public, to bear witness\textquotedblright{} (EtymDB, Lewis \& Short), built from \emph{prō} \textquotedblleft{}in public\textquotedblright{} and \emph{testārī} \textquotedblleft{}to bear witness\textquotedblright{}, on \emph{testis} \textquotedblleft{}witness\textquotedblright{}.
\end{cousinweb}

\begin{grammarlens}[title={What you've built}]
One more word for telling, explaining and describing at B1.
\end{grammarlens}

\subsection*{Your turn}
\begin{itemize}
\raggedright
  \item \emph{Say it:} \emph{protestar}
  \item \emph{Say it:} \emph{protestar}, once more
  \item \emph{Say it:} a sentence of your own with \emph{protestar}
  \item \emph{Recall:} say the Portuguese for the freedom, then the Portuguese for the equality, then say \emph{protestar} again
\end{itemize}

\begin{tcolorbox}[breakable,colback=teal!4,colframe=teal!35!black,title={Before you move on}]
What does protestar mean? (\textquotedblleft{}To protest\textquotedblright{}.) Say it once more.
\end{tcolorbox}
\section[o protesto]{o protesto — the protest}

\label{lesson:PT-C364-protesto}



\begin{quote}
Before the new one: say the Portuguese for the equality, then the Portuguese for to protest.
\end{quote}

\subsection*{You'll want to know: o protesto}
\textbf{o protesto} — \textquotedblleft{}the protest\textquotedblright{}.

\begin{quote}
Houve um protesto em frente à câmara.
\end{quote}

\begin{quote}
\emph{There was a protest in front of the town hall.}
\end{quote}

\emph{Em protesto contra} — in protest against.

\begin{cousinweb}
\emph{protesto} is built on \emph{protestar} \textquotedblleft{}to protest\textquotedblright{}, from Latin \emph{prōtestārī} \textquotedblleft{}to declare in public, to bear witness\textquotedblright{}, on \emph{testis} \textquotedblleft{}witness\textquotedblright{} (Lewis \& Short).
\end{cousinweb}

\begin{grammarlens}[title={What you've built}]
One more word for telling, explaining and describing at B1.
\end{grammarlens}

\subsection*{Your turn}
\begin{itemize}
\raggedright
  \item \emph{Say it:} \emph{o protesto}
  \item \emph{Say it:} \emph{o protesto}, once more
  \item \emph{Say it:} a sentence of your own with \emph{o protesto}
  \item \emph{Recall:} say the Portuguese for the equality, then the Portuguese for to protest, then say \emph{o protesto} again
\end{itemize}

\begin{tcolorbox}[breakable,colback=teal!4,colframe=teal!35!black,title={Before you move on}]
What does o protesto mean? (\textquotedblleft{}The protest\textquotedblright{}.) Say it once more.
\end{tcolorbox}
\section[a manifestação]{a manifestação — the demonstration, the march; the expression (of a feeling)}

\label{lesson:PT-C364-manifestacao}



\begin{quote}
Before the new one: say the Portuguese for to protest, then the Portuguese for the protest.
\end{quote}

\subsection*{You'll want to know: a manifestação}
\textbf{a manifestação} — \textquotedblleft{}the demonstration, the march; the expression (of a feeling)\textquotedblright{}.

\begin{quote}
Milhares de pessoas foram à manifestação.
\end{quote}

\begin{quote}
\emph{Thousands of people went to the demonstration.}
\end{quote}

In speech \emph{a manif}.

\begin{cousinweb}
\emph{manifestação} is a learned borrowing from Latin \emph{manifēstātiō} \textquotedblleft{}a manifesting, a manifestation\textquotedblright{} (EtymDB, Lewis \& Short), from \emph{manifēstāre} \textquotedblleft{}to show\textquotedblright{}, on \emph{manifēstus} \textquotedblleft{}clear, evident\textquotedblright{}.
\end{cousinweb}

\begin{grammarlens}[title={What you've built}]
One more word for telling, explaining and describing at B1.
\end{grammarlens}

\subsection*{Your turn}
\begin{itemize}
\raggedright
  \item \emph{Say it:} \emph{a manifestação}
  \item \emph{Say it:} \emph{a manifestação}, once more
  \item \emph{Say it:} a sentence of your own with \emph{a manifestação}
  \item \emph{Recall:} say the Portuguese for to protest, then the Portuguese for the protest, then say \emph{a manifestação} again
\end{itemize}

\begin{tcolorbox}[breakable,colback=teal!4,colframe=teal!35!black,title={Before you move on}]
What does a manifestação mean? (\textquotedblleft{}The demonstration, the march; the expression (of a feeling)\textquotedblright{}.) Say it once more.
\end{tcolorbox}
\section[manifestar]{manifestar — to show, to express (a feeling); (manifestar-se) to demonstrate}

\label{lesson:PT-C364-manifestar}



\begin{quote}
Before the new one: say the Portuguese for the protest, then the Portuguese for the demonstration.
\end{quote}

\subsection*{You'll want to know: manifestar}
\textbf{manifestar} — \textquotedblleft{}to show, to express (a feeling); (manifestar-se) to demonstrate\textquotedblright{}.

\begin{quote}
Os estudantes manifestaram-se no centro da cidade.
\end{quote}

\begin{quote}
\emph{The students demonstrated in the city centre.}
\end{quote}

\emph{Manifestar-se} — to demonstrate; \emph{manifestar uma opinião} — to express an opinion.

\begin{cousinweb}
\emph{manifestar} is a learned borrowing from Latin \emph{manifēstāre} \textquotedblleft{}to show, to make public\textquotedblright{} (EtymDB), on \emph{manifēstus} \textquotedblleft{}clear, evident\textquotedblright{} (Lewis \& Short).
\end{cousinweb}

\begin{grammarlens}[title={What you've built}]
One more word for telling, explaining and describing at B1.
\end{grammarlens}

\subsection*{Your turn}
\begin{itemize}
\raggedright
  \item \emph{Say it:} \emph{manifestar}
  \item \emph{Say it:} \emph{manifestar}, once more
  \item \emph{Say it:} a sentence of your own with \emph{manifestar}
  \item \emph{Recall:} say the Portuguese for the protest, then the Portuguese for the demonstration, then say \emph{manifestar} again
\end{itemize}

\begin{tcolorbox}[breakable,colback=teal!4,colframe=teal!35!black,title={Before you move on}]
What does manifestar mean? (\textquotedblleft{}To show, to express (a feeling); (manifestar-se) to demonstrate\textquotedblright{}.) Say it once more.
\end{tcolorbox}
\section[exigir]{exigir — to demand, to require}

\label{lesson:PT-C364-exigir}



\begin{quote}
Before the new one: say the Portuguese for the demonstration, then the Portuguese for to show.
\end{quote}

\subsection*{You'll want to know: exigir}
\textbf{exigir} — \textquotedblleft{}to demand, to require\textquotedblright{}.

\begin{quote}
Os pais exigem mais segurança na escola.
\end{quote}

\begin{quote}
\emph{The parents are demanding more safety at the school.}
\end{quote}

\emph{Exigir que} takes the subjunctive.

\begin{cousinweb}
\emph{exigir} is a learned borrowing from Latin \emph{exigere} \textquotedblleft{}to drive out, to demand\textquotedblright{} (EtymDB), built from \emph{ex} and \emph{agere} \textquotedblleft{}to drive\textquotedblright{} (Lewis \& Short).
\end{cousinweb}

\begin{grammarlens}[title={What you've built}]
One more word for telling, explaining and describing at B1.
\end{grammarlens}

\subsection*{Your turn}
\begin{itemize}
\raggedright
  \item \emph{Say it:} \emph{exigir}
  \item \emph{Say it:} \emph{exigir}, once more
  \item \emph{Say it:} a sentence of your own with \emph{exigir}
  \item \emph{Recall:} say the Portuguese for the demonstration, then the Portuguese for to show, then say \emph{exigir} again
\end{itemize}

\begin{tcolorbox}[breakable,colback=teal!4,colframe=teal!35!black,title={Before you move on}]
What does exigir mean? (\textquotedblleft{}To demand, to require\textquotedblright{}.) Say it once more.
\end{tcolorbox}
